\documentclass[10pt,a4paper]{amsart}
\usepackage{amsmath,amssymb,amsthm,graphicx}
% Retained preamble for the Stacks Project chapter divisors.tex (Divisors).
% The chapter file pinned-original-source.tex carries only \input{preamble} on its first line, so no
% preamble text lives inside the chapter itself.  The Stacks Project uses one single shared preamble for
% every chapter; the snapshot stages it verbatim in sources/stacks-chow/preamble.tex.  The block below is
% copied mechanically, line for line, from that file's lines 14-18 (the Xy-pic setup) and lines 154-217
% (the native theorem declarations together with the mathematical macros).  Nothing was added, removed or
% redefined.  The wrapper class and the amsmath/amssymb/amsthm/graphicx base packages are supplied by the
% assembler, so the class/\externaldocument part of the shared preamble is not repeated here.
% For commutative diagrams we use Xy-pic
\usepackage[all]{xy}

% We use 2cell for 2-commutative diagrams.
\xyoption{2cell}
% Theorem environments.
%
\theoremstyle{plain}
\newtheorem{theorem}[subsection]{Theorem}
\newtheorem{proposition}[subsection]{Proposition}
\newtheorem{lemma}[subsection]{Lemma}

\theoremstyle{definition}
\newtheorem{definition}[subsection]{Definition}
\newtheorem{example}[subsection]{Example}
\newtheorem{exercise}[subsection]{Exercise}
\newtheorem{situation}[subsection]{Situation}

\theoremstyle{remark}
\newtheorem{remark}[subsection]{Remark}
\newtheorem{remarks}[subsection]{Remarks}

\numberwithin{equation}{subsection}

% Macros
%
\def\lim{\mathop{\mathrm{lim}}\nolimits}
\def\colim{\mathop{\mathrm{colim}}\nolimits}
\def\Spec{\mathop{\mathrm{Spec}}}
\def\Hom{\mathop{\mathrm{Hom}}\nolimits}
\def\Ext{\mathop{\mathrm{Ext}}\nolimits}
\def\SheafHom{\mathop{\mathcal{H}\!\mathit{om}}\nolimits}
\def\SheafExt{\mathop{\mathcal{E}\!\mathit{xt}}\nolimits}
\def\Sch{\mathit{Sch}}
\def\Mor{\mathop{\mathrm{Mor}}\nolimits}
\def\Ob{\mathop{\mathrm{Ob}}\nolimits}
\def\Sh{\mathop{\mathit{Sh}}\nolimits}
\def\NL{\mathop{N\!L}\nolimits}
\def\CH{\mathop{\mathrm{CH}}\nolimits}
\def\proetale{{pro\text{-}\acute{e}tale}}
\def\etale{{\acute{e}tale}}
\def\QCoh{\mathit{QCoh}}
\def\Ker{\mathop{\mathrm{Ker}}}
\def\Im{\mathop{\mathrm{Im}}}
\def\Coker{\mathop{\mathrm{Coker}}}
\def\Coim{\mathop{\mathrm{Coim}}}

% Boxtimes
%
\DeclareMathSymbol{\boxtimes}{\mathbin}{AMSa}{"02}

%
% Macros for moduli stacks/spaces
%
\def\QCohstack{\mathcal{QC}\!\mathit{oh}}
\def\Cohstack{\mathcal{C}\!\mathit{oh}}
\def\Spacesstack{\mathcal{S}\!\mathit{paces}}
\def\Quotfunctor{\mathrm{Quot}}
\def\Hilbfunctor{\mathrm{Hilb}}
\def\Curvesstack{\mathcal{C}\!\mathit{urves}}
\def\Polarizedstack{\mathcal{P}\!\mathit{olarized}}
\def\Complexesstack{\mathcal{C}\!\mathit{omplexes}}
% \Pic is the operator that assigns to X its picard group, usage \Pic(X)
% \Picardstack_{X/B} denotes the Picard stack of X over B
% \Picardfunctor_{X/B} denotes the Picard functor of X over B
\def\Pic{\mathop{\mathrm{Pic}}\nolimits}
\def\Picardstack{\mathcal{P}\!\mathit{ic}}
\def\Picardfunctor{\mathrm{Pic}}
\def\Deformationcategory{\mathcal{D}\!\mathit{ef}}

% --- appended 2026-10-12 (Chief): the project own definitions of the reference/slogan/history
% display environments, copied mechanically from the shared Stacks preamble lines 8-12
% (\usepackage{verbatim} plus the three \newenvironment lines, which the chapter preamble keeps
% verbatim from the project preamble). Nothing added, removed or redefined: these environments
% carry no mathematical content and the project itself typesets them as comments. Without these
% three lines a counted statement whose own text uses them fails with Environment slogan undefined.
\usepackage{verbatim}
\newenvironment{reference}{\comment}{\endcomment}
\newenvironment{slogan}{\comment}{\endcomment}
\newenvironment{history}{\comment}{\endcomment}
\usepackage{cleveref}
\crefname{originallemma}{original lemma}{original lemmas}
\Crefname{originallemma}{Original lemma}{Original lemmas}
\makeatletter
\newlabel{properties-lemma-quasi-affine-invertible-nonvanishing-section}{{properties-lemma-quasi-affine-invertible-nonvanishing-section}{source}{}{sourcepointer-1}{}}
\newlabel{properties-lemma-quasi-affine-invertible-nonvanishing-section@cref}{{[originallemma][0][]properties-lemma-quasi-affine-invertible-nonvanishing-section}{[1][0][]source}{}{}{}}
\makeatother
\crefname{originallemma}{original lemma}{original lemmas}
\Crefname{originallemma}{Original lemma}{Original lemmas}
\makeatletter
\newlabel{schemes-lemma-equivalence-quasi-coherent}{{schemes-lemma-equivalence-quasi-coherent}{source}{}{sourcepointer-2}{}}
\newlabel{schemes-lemma-equivalence-quasi-coherent@cref}{{[originallemma][0][]schemes-lemma-equivalence-quasi-coherent}{[1][0][]source}{}{}{}}
\makeatother
\title{An invertible module on a finite cover is trivial over the preimage of a neighbourhood of any point of the base}
\author{The Stacks Project Authors}
\date{stacks-project:divisors.tex:local-snapshot-20261012; GNU FDL 1.2 or later}
\begin{document}
\maketitle
\noindent\textit{Editorial note.} All delivered mathematical text is the author's own, copied from the pinned Stacks Project chapter file pinned-original-source.tex; no statement or proof step was rewritten, reordered or removed.\par\medskip
\noindent\textit{Editorial note.} Following the multi-claim convention for this source, source\_url is the source-level URL https://stacks.math.columbia.edu/ and the source label lemma-finite-trivialize-invertible-upstairs is carried in claim\_key and here rather than in a per-claim tag URL.\par\medskip
\noindent\textit{Editorial note.} The preamble is the chapter's shared Stacks Project preamble, retained as sources/stacks-divisors/retained-preamble.tex: a line-for-line excerpt of the shared Stacks preamble staged with this snapshot. No macro was invented or redefined, and this card adds no preamble text of its own.\par\medskip
\noindent\textit{Editorial note.} Statement span 3616-3622 and proof span 3624-3652 come from the halfA pool card operations/pool-cards/stacks-divisors-halfA.json and were re-read in the pinned chapter file before this card was written.\par\medskip
\noindent\textit{Editorial note.} The proof contains the author's sentence offering a cited alternative ('This follows from Properties, Lemma ... but we also give a direct argument.'); both halves are preserved and the direct argument is the one counted.\par\medskip
\noindent\textit{Editorial note.} 2 cross-references occur inside the counted spans; they are listed in external\_references as navigation pointers only.\par\medskip
\noindent\textit{Allowed standard prerequisites (editorial).}\begin{itemize}
\item A finite morphism is closed and has finite fibres, and the fibre of a finite morphism over a point of an affine base is finite. Used at proof lines 3625-3626 and at 3649-3651. Gives both the finiteness of the fibre used to select the auxiliary set of closed points and the closedness used to push the open set down.
\item Surjectivity of a map of quasi-coherent modules over an affine scheme after taking global sections, and the identification of global sections with the sum of the fibres at a finite set meeting every specialization class. Used at proof lines 3633-3643 with the cited schemes chapter lemma on the equivalence of quasi-coherent modules. Produces a global section of the invertible module that does not vanish at any point of the fibre.
\item Existence of a section of an invertible module on an affine scheme not vanishing at a prescribed finite set of points, cited from the properties chapter as the alternative to the direct argument given here. Cited at proof lines 3628-3630. The author offers it as the shortcut and then gives the direct argument that is counted here.
\end{itemize}\medskip
\ifdefined\hypertarget\hypertarget{sourcepointer-1}{}\fi
\noindent\textit{Source reference.} Original lemma properties-lemma-quasi-affine-invertible-nonvanishing-section: sources/stacks-divisors/pinned-original-source.tex, source lines 3241--3241. Same-chapter reference used in the retained author text at source line(s) 3241, 3630. It is supplied only as a navigation pointer so that the retained text compiles; the counted item delivers the author's own statement and the author's complete proof, and no part of the referenced item is copied in or substituted for the counted argument.\par\medskip
\ifdefined\hypertarget\hypertarget{sourcepointer-2}{}\fi
\noindent\textit{Source reference.} Original lemma schemes-lemma-equivalence-quasi-coherent: sources/stacks-divisors/pinned-original-source.tex, source lines 3645--3645. Same-chapter reference used in the retained author text at source line(s) 3645. It is supplied only as a navigation pointer so that the retained text compiles; the counted item delivers the author's own statement and the author's complete proof, and no part of the referenced item is copied in or substituted for the counted argument.\par\medskip
\begin{lemma}
\label{lemma-finite-trivialize-invertible-upstairs}
Let $\pi : X \to Y$ be a finite morphism of schemes.
Let $\mathcal{L}$ be an invertible $\mathcal{O}_X$-module.
Let $y \in Y$. There exists an open neighbourhood
$V \subset Y$ of $y$ such that $\mathcal{L}|_{\pi^{-1}(V)}$ is trivial.
\end{lemma}

\begin{proof}
Clearly we may assume $Y$ and hence $X$ affine. Since $\pi$ is finite the
fibre $\pi^{-1}(\{y\})$ over $y$ is finite.
Since $X$ is affine, we can pick $s \in \Gamma(X, \mathcal{L})$
not vanishing in any point of $\pi^{-1}(\{y\})$. This follows
from Properties, Lemma
\ref{properties-lemma-quasi-affine-invertible-nonvanishing-section}
but we also give a direct argument. Namely, we can
pick a finite set $E \subset X$ of closed points such that
every $x \in \pi^{-1}(\{y\})$ specializes to some point of $E$.
For $x \in E$ denote $i_x : x \to X$ the closed immersion.
Then
$\mathcal{L} \to \bigoplus_{x \in E} i_{x, *}i_x^*\mathcal{L}$
is a surjective map of quasi-coherent $\mathcal{O}_X$-modules,
and hence the map
$$
\Gamma(X, \mathcal{L}) \to
\bigoplus\nolimits_{x \in E} \mathcal{L}_x/\mathfrak m_x\mathcal{L}_x
$$
is surjective (as taking global sections is an exact functor on the
category of quasi-coherent $\mathcal{O}_X$-modules, see
Schemes, Lemma \ref{schemes-lemma-equivalence-quasi-coherent}).
Thus we can find an $s \in \Gamma(X, \mathcal{L})$
not vanishing at any point specializing to a point of $E$.
Then $X_s \subset X$ is an open neighbourhood of $\pi^{-1}(\{y\})$.
Since $\pi$ is finite, hence closed, we conclude that there is an
open neighbourhood $V \subset Y$ of $y$ whose inverse image
is contained in $X_s$ as desired.
\end{proof}

\end{document}
